%% GENERIC STYLE SETTINGS
\usetheme{default}
\usepackage{listings}
\usepackage{multirow}
\usepackage{xcolor}

\usebackgroundtemplate{
\includegraphics[width=\paperwidth,height=\paperheight]{images/hutton_background}
}
%% PRESENTATION CONFIGURATION PARAMETERS %%%%%%%%%%%%%%%%%%%%%%%%%%%%%%%%%%%%%%%
%\titlebackgroundfile{images/hutton_title}
%\framebackgroundfile{images/hutton_background}
\definecolor{hutton_green}{HTML}{78A22F}
\definecolor{hutton_purple}{HTML}{872175}
\definecolor{hutton_blue}{HTML}{569BBE}
\usefonttheme{structurebold}
\setbeamercolor{alerted text}{fg=orange}
\setbeamercolor{background canvas}{bg=white}
\setbeamercolor{block title}{bg=hutton_purple}
\setbeamercolor{frametitle}{fg=hutton_purple}
\setbeamercolor{title}{fg=black}
\setbeamercolor{titlelike}{fg=hutton_green}
\setbeamercolor{author}{fg=hutton_purple}
\setbeamercolor{author in head/foot}{fg=white}
\setbeamercolor{title in head/foot}{fg=white}
\setbeamercolor{section in head/foot}{fg=hutton_purple}
\setbeamercolor{normal text}{fg=black}
\setbeamercolor{frametitle}{fg=hutton_purple}
\setbeamerfont{block title}{size={}}
\setbeamerfont{author}{size=\footnotesize}
\setbeamerfont{date}{size=\footnotesize}
\setbeamercolor{section in toc shaded}{fg=hutton_purple}
\setbeamercolor{section in toc}{fg=hutton_purple}
\setbeamercolor{subsection in toc shaded}{fg=hutton_purple}
\setbeamercolor{subsection in toc}{fg=hutton_purple}
\setbeamertemplate{itemize item}[circle]
\setbeamertemplate{itemize subitem}[circle]
\setbeamertemplate{itemize subsubitem}[circle]
\setbeamertemplate{itemize subsubsubitem}[circle]
\setbeamercolor{itemize item}{fg=hutton_purple}
\setbeamercolor{itemize subitem}{fg=hutton_purple}
\setbeamercolor{itemize subsubitem}{fg=hutton_purple}
\setbeamercolor{itemize subsubsubitem}{fg=hutton_purple}
\setbeamercolor{enumerate item}{fg=hutton_purple}
\setbeamercolor{enumerate subitem}{fg=hutton_purple}
\setbeamercolor{enumerate subsubitem}{fg=hutton_purple}
\setbeamercolor{enumerate subsubsubitem}{fg=hutton_purple}
\setbeamercolor{alerted text}{fg=hutton_green}
\setbeamerfont{alerted text}{series=\bfseries}
% This command makes sure that acrobat reader doesn't change the colours of the slide
% when there are figures with transparencies.
\pdfpageattr {/Group << /S /Transparency /I true /CS /DeviceRGB>>}

%Disables discrete bottom navigation bar
%\beamertemplatenavigationsymbolsempty

% Modify the slide titles to avoid the corner images,
\setbeamertemplate{frametitle}
{
\vspace{0.05\textheight}
\noindent\quad\begin{minipage}[t][0.12\textheight][t]{0.85\textwidth}
\insertframetitle\par
\end{minipage}
}

% Modify title page to avoid the big logo on right
\setbeamertemplate{title page}{
    \begin{picture}(0,0)
            %This ends up on top of the default background image, rather than replacing it:
            \put(-30,-165){%
                \includegraphics[width=\paperwidth,height=\paperheight]{images/hutton_title}
            }
            \put(0,-75){%
                \begin{minipage}[b][0.4\textheight][t]{0.75\textwidth}
                    \usebeamerfont{title}\usebeamercolor[fg]{title}{\inserttitle\par}
                    \usebeamerfont{subtitle}\usebeamercolor[fg]{subtitle}{\insertsubtitle\par}
                \end{minipage}
            }
            \put(0,-135){%
                \begin{minipage}[b][0.1\textheight][t]{\textwidth}
                    \usebeamerfont{author}\usebeamercolor[fg]{author}{\insertauthor\par}
                \end{minipage}
            }
    \end{picture}
}

%%%%%%%%%%%%%%%%%%%%%%%%%%%%%%%%%%%%%%%%%%%%%%%%%%%%%%%%%%%%%%%%%%%%%%%%%%%%%%%%

% LISTINGS SETTING
% Settings for code listings in lstlistings

\definecolor{hutton_lightgreen}{HTML}{C8F27F}

\lstset{ %
  backgroundcolor=\color{hutton_lightgreen},   % choose the background color; you must add \usepackage{color} or \usepackage{xcolor}
  basicstyle=\tiny\ttfamily,        % the size of the fonts that are used for the code
  breakatwhitespace=true,          % sets if automatic breaks should only happen at whitespace
  breaklines=true,                 % sets automatic line breaking
  captionpos=b,                    % sets the caption-position to bottom
  commentstyle=\color{red},    % comment style
  deletekeywords={...},            % if you want to delete keywords from the given language
  escapeinside={\%*}{*)},          % if you want to add LaTeX within your code
  extendedchars=true,              % lets you use non-ASCII characters; for 8-bits encodings only, does not work with UTF-8
  frame=single,                    % adds a frame around the code
  keepspaces=true,                 % keeps spaces in text, useful for keeping indentation of code (possibly needs columns=flexible)
  keywordstyle=\color{blue},       % keyword style
%  language=Octave,                 % the language of the code
  morekeywords={*,...},            % if you want to add more keywords to the set
  numbers=left,                    % where to put the line-numbers; possible values are (none, left, right)
  numbersep=5pt,                   % how far the line-numbers are from the code
  numberstyle=\tiny\color{gray}, % the style that is used for the line-numbers
  rulecolor=\color{black},         % if not set, the frame-color may be changed on line-breaks within not-black text (e.g. comments (green here))
  showspaces=false,                % show spaces everywhere adding particular underscores; it overrides 'showstringspaces'
  showstringspaces=false,          % underline spaces within strings only
  showtabs=false,                  % show tabs within strings adding particular underscores
  stepnumber=1,                    % the step between two line-numbers. If it's 1, each line will be numbered
  stringstyle=\color{violet},     % string literal style
  tabsize=4,                       % sets default tabsize to 2 spaces
  title=\lstname                   % show the filename of files included with \lstinputlisting; also try caption instead of title
}

% Misc packages
\usepackage{multicol}

%%%
% TITLE PREAMBLE
\title[Biopython Project Update 2018] % (optional, only for long titles)
{Biopython Project Update 2018}
\subtitle{Standing on each other's shoulders \\ \includegraphics[height=3cm]{images/biopython_logo_m.png}}
\author[Cock] % (optional, for multiple authors)
{Peter~Cock (\href{https://twitter.com/pjacock}{@pjacock on Twitter}), \\
The Biopython Contributors (\href{https://twitter.com/Biopython}{@biopython on Twitter})}
\institute[The James Hutton Institute] % (optional)
{
  Information and Computational Sciences\\
  The James Hutton Institute
}
\date[June 2018] % (optional)
{Portland, Oregon, USA, 28$^{th}$ June 2018}
\subject{Bioinformatics}

%%%
% START DOCUMENT
\begin{document}

\frame[plain]{\titlepage}

\begin{frame}
  \frametitle{What is Biopython?}

  \begin{itemize}
  \item Collection of modules for biological computation in Python
  \begin{itemize}
  \item Sequence handling and motifs, parsers, database queries, protein structures, phylogenetics, tool wrappers and more.
  \end{itemize}
  \item Started in 1999, first release in 2000
  \item Open source and freely available (Biopython license)
  \item \url{https://biopython.org} and @Biopython on Twitter
  \end{itemize}

\center
\includegraphics[height=2.5cm]{images/biopython_logo_m.png}
\end{frame}


\begin{frame}
  \frametitle{43 named contributors in last year, \\ 18 newcomers with star!}
  \scriptsize{
  \begin{multicols}{3}
  \begin{itemize}
    \item Adhemar Zerlotini
    \item Ariel Aptekmann
    \item Benjamin Vaisvil*
    \item Blaise Li
    \item Chad Parmet
    \item Chris Rands
    \item Christian Brueffer
    \item Connor T. Skennerton
    \item Erik Cederstrand*
    \item Fei Qi*
    \item Francesco Gastaldello
    \item James Jeffryes*
    \item Jerven Bolleman*
    \item Joe Greener*
    \item Joerg Schaarschmidt*
    \item João Rodrigues
    \item Jeroen Van Goey
    \item Jun Aruga*
    \item Kai Blin
    \item Kozo Nishida
    \item Lewis A. Marshall*
    \item Markus Piotrowski
    \item Michiel de Hoon
    \item Nicolas Fontrodona*
    \item Pamela Russell*
    \item Peter Cock
    \item Philip Bergstrom*
    \item rht*
    \item Saket Choudhary
    \item Shuichiro Makigaki*
    \item Shyam Saladi*
    \item Siong Kong
    \item Spencer Bliven
    \item Stefans Mezulis
    \item Steve Bond
    \item Yasar L. Ahmed*
    \item Zachary Sailer*
    \item Zaid Ur-Rehman*
    \item Wibowo 'Bow' Arindrarto
  \end{itemize}
  \end{multicols}
  }
\end{frame}


\begin{frame}
\frametitle{Biopython 1.71 (April 2018)}
\begin{itemize}
   \item Python 3-compatible rich comparison for all Bio.PDB objects
   \item New options and fixes for parsing and writing files, including PIR format for sequences, mmCIF files, and Gene files.
   \item New codon tables 27-31 from NCBI (NCBI genetic code table version 4.2) were added.
       \begin{itemize}
          \item Now supports ambiguous codons for either stop or amino acid
       \end{itemize}
\end{itemize}
\end{frame}

\begin{frame}
\frametitle{Biopython 1.72 (June 2018)}
\begin{itemize}
    \item NumPy no longer needed at compile time
    \item Bio.Entrez now supports NCBI API key
    \item New alternative pairwise alignment module
    \item Bio.SearchIO is no longer considered experimental
\end{itemize}
\end{frame}

\begin{frame}
\frametitle{Python Versions}
\begin{itemize}
\item Currently support Python 2.7, 3.4, 3.5 and 3.6
\item Clear end of life for Python 2 support in 2020, pledged on \url{http://python3statement.org/}
\item Also support and test on PyPy
\item Support for Jython now deprecated
\item Next release will officially support Python 3.7
\end{itemize}
\end{frame}


\begin{frame}
\frametitle{PyPI, pip and wheels}
\begin{itemize}
\item Recommend ``pip install biopython``
\item This fetches from Python Package Index (PyPI)
\item Will use pre-compiled wheel files if available
\item We build the wheels on AppVeyor \& TravisCI
    \begin{itemize}
        \item Following NumPy community in using the multibuild system, \\
              developed by Matthew Brett and the MacPython project. \\
              \url{https://github.com/matthew-brett/multibuild} \\
              \url{https://github.com/biopython/biopython-wheels}
    \end{itemize}
\end{itemize}
\end{frame}


\begin{frame}
\frametitle{Biopython's Open Source License}
\begin{itemize}
\item The Open Source Initiative \url{https://opensource.org/}
    maintains a list of approved open source licenses
\item This is the \textit{de facto} gold standard.
\item Biopython's license is not (quite) on that list.
\end{itemize}
\end{frame}

\begin{frame}
\frametitle{Historical Permission Notice \\ and Disclaimer (HPND)}
Permission to use, copy, modify and distribute this software and its documentation for any purpose and without fee is hereby granted, provided that the above copyright notice appear in all copies[,] [and] that both [that] [the] copyright notice and this permission notice appear in supporting documentation[, and that the name [of] $<$copyright holder$>$ [or $<$related entities$>$] not be used in advertising or publicity pertaining to distribution of the software without specific, written prior permission]. [$<$copyright holder$>$ makes no representations about the suitability of this software for any purpose. It is provided ``as is'' without express or implied warranty.]

%[<copyright holder> DISCLAIMS ALL WARRANTIES WITH REGARD TO THIS SOFTWARE, INCLUDING ALL IMPLIED WARRANTIES OF MERCHANTABILITY AND FITNESS[,][.] IN NO EVENT SHALL <copyright holder> BE LIABLE FOR ANY SPECIAL, INDIRECT OR CONSEQUENTIAL DAMAGES OR ANY DAMAGES WHATSOEVER RESULTING FROM LOSS OF USE, DATA OR PROFITS, WHETHER IN AN ACTION OF CONTRACT, NEGLIGENCE OR OTHER TORTIOUS ACTION, ARISING OUT OF OR IN CONNECTION WITH THE USE OR PERFORMANCE OF THIS SOFTWARE.]

\url{https://opensource.org/licenses/HPND}
\end{frame}

\begin{frame}
\frametitle{Biopython License Agreement}
\textit{Italics highlighting ony substantial difference to the HPND:}
\vspace{6pt} \\

Permission to use, copy, modify, and distribute this software and its documentation \textit{with or without modifications} and for any purpose and without fee is hereby granted, provided that any copyright notices appear in all copies and that both those copyright notices and this permission notice appear in supporting documentation, and that the names of the contributors or copyright holders not be used in advertising or publicity pertaining to distribution of the software without specific prior permission.

%THE CONTRIBUTORS AND COPYRIGHT HOLDERS OF THIS SOFTWARE DISCLAIM ALL WARRANTIES WITH REGARD TO THIS SOFTWARE, INCLUDING ALL IMPLIED WARRANTIES OF MERCHANTABILITY AND FITNESS, IN NO EVENT SHALL THE CONTRIBUTORS OR COPYRIGHT HOLDERS BE LIABLE FOR ANY SPECIAL, INDIRECT OR CONSEQUENTIAL DAMAGES OR ANY DAMAGES WHATSOEVER RESULTING FROM LOSS OF USE, DATA OR PROFITS, WHETHER IN AN ACTION OF CONTRACT, NEGLIGENCE OR OTHER TORTIOUS ACTION, ARISING OUT OF OR IN CONNECTION WITH THE USE OR PERFORMANCE OF THIS SOFTWARE.
\end{frame}

\begin{frame}
\frametitle{Biopython's Open Source License}
\begin{itemize}
\item We could ask the Open Source Initiative (OSI)
    to approve our license (as an HPND variant?)
\item We've agreed to gradually dual-license the code
    under our old license and the 3-clause BSD license
\item Requires checking each file to confirm all
    contributors agree to dual licensing
\item As of Biopython 1.72, about $20\%$ of code files done
\item Only one minor contributor has refused to date
\end{itemize}
\end{frame}


\begin{frame}
\frametitle{On going work}
\begin{itemize}
    \item Moving from epydoc to Sphinx for API docs \\
          (and from \LaTeX to Sphinx for Tutorial?)
    \item Improving compliance with PEP8 and PEP257 style guidelines
    \item Improving code test coverage \url{https://codecov.io/github/biopython/biopython/}
    \item Planning what to do with legacy Alphabet objects
    \item Simplifying our release process
    \item Other new contributor driven efforts
\end{itemize}
\end{frame}


\begin{frame}
\frametitle{Changes to help Community Building}
\begin{itemize}
   \item GitHub Issue templates (to help with bug reporting)
   \item Pull request templates (to help with expectations)
   \item \texttt{CODEOWNERS} file to help assign code reviews
   \item Using \textit{Easy Fix} tag on some issues, \\ intended for new contributors
   \item Already have a \texttt{CONTRIBUTING} text file
   \item Do we need a \texttt{CODE\_OF\_CONDUCT} too?
   \item \textbf{What else should we be doing?}
\end{itemize}
\end{frame}


\begin{frame}
\frametitle{Acknowledgements}
Thank you to:
\begin{itemize}
    \item All our contributors to date
    \item Google Summer of Code (supporting past students)
    \item My travel costs paid for by:
    \begin{itemize}
       \item James Hutton Institute
       \item Open Bioinformatics Foundation \\
             (as a BOSC committee member)
    \end{itemize}
\end{itemize}

\vspace{1.1cm}
\center\includegraphics[width=0.9\paperwidth]{../images/Hutton-thanks-banner}
\end{frame}

% etc
\end{document}
